\documentclass[11pt]{article}
\usepackage{amsmath,amssymb,amsthm,mathtools,hyperref}

\begin{document}
\section*{Origin of self-replicating ribozymes}
\label{prob:287}
\label{rank:190}
\par\smallskip

\noindent\textbf{Domain.} RNA replication; autonomous propagation. \quad \textbf{Retained status/source context.} Open sustained, chemically integrated propagation question, explicitly narrowed after 2026 self/complement synthesis and 2025 exponential RNA replication advances.

\paragraph{Definitions and assumptions.} Let $R$ be a finite set of functional RNA lineage types, defined by catalyst sequences or compartmentalized RNA assemblages. Count reproducing units $N_i$ under a repeated environment and serial-dilution protocol. Feedstocks are short substrates with specified activation chemistry; modern proteins and external catalyst replacement are excluded after initialization. Each concrete instance declares substrate lengths and activation states, supplied sequence-bearing materials, other material inputs, externally imposed cycling and the resource/energy boundary. The admissible range of ``short'' substrates and of externally supplied sequence information is not fixed here. For the branching variant assume independent offspring, finite offspring means, identical cycle conditions, an irreducible nonsingular offspring law and negligible between-lineage competition. Independence here concerns the offspring vectors of distinct reproducing units, conditional on their types, under a time-invariant type-specific law; components within one offspring vector need not be independent. The chosen unit and type assignment must support that law, with complementary strands, catalytic partners, mutants outside $R$ and nonfunctional products explicitly accounted for. Then $K_{ij}$ is expected viable type-$i$ offspring per type-$j$ unit after copying, strand separation, decay and dilution. Assays distinguish renewed catalysts from retained seeds. Cycle-boundary counts and assignment of offspring to starting units must be declared. Whether surviving parental units enter $K$ remains unspecified; separate measurement of renewal does not determine this transition-counting convention.

\paragraph{Problem statement.} Find experimentally supported conditions with
\[
\rho(K)>1
\]
and independently verify functional propagation and heritable selectable variation over $K_0$ cycles. Here $\rho$ is spectral radius and $K_0$ is declared in advance. For each concrete instance, declare the positive integer $K_0$, cycle duration, functional phenotype, renewal and inheritance readouts, and the comparison supporting selectable variation, together with the observation law, uncertainty assessment and success criteria required by the binding catalogue conventions. No particular assay or numerical threshold is selected here. Under the stated stationary branching assumptions, this gives positive survival probability, not certain survival. Then connect substrate replenishment to a chemically plausible energy source. A positive single-pass synthesis yield is insufficient to establish this net reproductive criterion.

\paragraph{Additional binding assumptions and resolution criterion.} The spectral-radius survival criterion is established for the stated multitype branching assumptions; the unresolved part is a chemically realized, measured, self-renewing offspring process that satisfies it. Specify its complete offspring generating law to quantify survival probability and initial-type effects. Distinguish the specified law and initial population law from quantities inferred from observations, and state the model restrictions and uncertainty assessment supporting the inference of $\rho(K)>1$. Finite-cycle observations do not identify an unrestricted complete offspring law or establish indefinite stationarity; indefinite survival probabilities are conditional predictions of the specified branching model. Keep cycle conditions and dilution fixed, track renewed catalyst rather than retained seed, and test the validity of negligible competition during the claimed propagation phase.

\paragraph{Alternative formulations.} (i) A supplied-substrate branching formulation tests catalyst renewal with independently reproducing units and replenished feeds. (ii) For cooperative density-dependent replication let $n_k$ be the abundance and resource state after cycle $k$, use a measured nonlinear map $n_{k+1}=F(n_k)$, and seek an attracting, positive invariant set separated from extinction for a declared inoculum range. The spectral criterion alone does not certify this second target. A stronger chemical-closure variant regenerates activated substrates within a constrained reaction network rather than supplies them externally.

\paragraph{Obstructions and partial progress.} Gianni and colleagues' QT45 is a 45-nucleotide polymerase ribozyme able to synthesize its complement and a copy of itself in reported assays; the abstract reports approximately 0.2 percent yields over 72 days. Attwater and colleagues demonstrate exponential replication of RNA strands and partial catalyst self-replication under cyclic conditions. These resolve important feasibility claims but do not alone establish the displayed complete functional reproduction cycle. The original statement is therefore retained only in this narrower, carefully measured form.

\paragraph{Selective external check in this revision.} SC09, SC10, SC11 (13 October 2026; access level and exact support are recorded in the source-check appendix). These selected checks are not an exhaustive present-day status search.

\paragraph{Sources (supplied provenance; not all rechecked).} \begin{itemize}
\item Supplied Wikipedia background; not used as scientific status verification. \url{https://en.wikipedia.org/wiki/RNA_world}
\item Malaterre, Jeancolas and Nghe, The Origin of Life: What Is the Question? (2022); conceptual framing, not problem-specific experimental verification. Metadata and indexed abstract were checked; direct page access was intermittent. \url{https://pubmed.ncbi.nlm.nih.gov/35594335/}
\item Gianni et al., A small polymerase ribozyme that can synthesize itself and its complementary strand (2026). \url{https://pubmed.ncbi.nlm.nih.gov/41678588/} (Primary QT45 result; single-pass synthesis has been demonstrated, while sustained autonomous propagation is a different target.)
\item Attwater et al., Trinucleotide substrates under pH-freeze-thaw cycles enable open-ended exponential RNA replication by a polymerase ribozyme (2025). \url{https://www.nature.com/articles/s41557-025-01830-y} (Primary RNA replication-cycle evidence, including partial ribozyme self-replication.)
\end{itemize}


\section*{Binding catalogue conventions}
This edition revises both supplied catalogues and preserves original problem numbers 1--330. It adds one globally unique label \verb|\label{rank:x}| per problem, where smaller $x$ means greater provisional difficulty. A rank label is an anchor at that problem's section: \verb|\ref{rank:x}| returns the section number, not the literal difficulty number. The visible \texttt{rank:x} text and the index give the difficulty number itself.

The entries are \emph{parameterized mathematical research schemes}, including empirical prediction, causal identification, historical reconstruction, theory and engineering. They are not 330 certified unsolved mathematical conjectures. A rigorous instance fixes the objects and restrictions below; changing those objects changes the question. Established identities, finite calculations, non-identification constructions and solved special cases are not being posed anew as open. Some residual tasks concern informative bounds, biological realization or external validation rather than a new theorem.

Every main definition and problem statement received a mathematical editorial review. Each entry now has an additional binding-conditions paragraph; direct corrections are separately recorded in the accompanying audit. The alternative formulations inherit the applicable conventions, but are different targets unless equivalence is proved. Where an old broad phrase conflicts with an explicit restriction below it, the explicit restriction controls.

\textbf{Literature boundary.} Thirteen selective primary-literature checks cover eighteen entries in this revision; their access levels and scope are recorded explicitly. The supplied bibliography and its earlier source-access descriptions are retained as provenance, not represented as newly verified. Neither this review nor a source's description proves the absence of later solutions. The date of this edition is 13 October 2026; there is no claim of exhaustive current-status certification for all 330 entries.

\subsection*{Binding conventions for a concrete instance}
\paragraph{Objects, quantifiers and non-vacuity.}
An instance declares measurable state, input, output and parameter spaces; a fixed candidate family; admissible initial laws, interventions and environments; a finite horizon or specified asymptotic regime; observation and sampling laws; and performance metrics, positive normalization scales and positive tolerances. These are chosen before fitting or testing, not adjusted after seeing failures. Shared parameters must genuinely be shared across the declared validation domain. Finite dimension alone does not restrict arbitrary real-number encodings: specify a finite library, bounded description length, or regularity/complexity bounds. Exclude constant or no-signal solutions through a fixed informative-output or baseline-improvement requirement where appropriate. A universally quantified result cannot be established by choosing one trivial favorable instance.

\paragraph{Probability laws and data.}
Every expectation or probability refers to the specified stochastic law, including initial uncertainty, process noise, observation noise and sampling as appropriate. A deterministic equation is not a probability model without a declared random input or initial law. ``Observed'' targets in population-level formulas mean population observables $g_j^*$, not noisy estimates asserted to be exact. Let $\widetilde g_j$ be independent validation estimates and suppose
\[
\mathbb{P}\{\lvert\widetilde g_j-g_j^*\rvert\le r_j\text{ for every }j\}\ge1-\alpha.
\]
Then an uncertainty-aware sufficient certification rule for positive target tolerances $\varepsilon_j$ is
\[
\max_j\frac{|\widehat g_j-\widetilde g_j|+r_j}{\varepsilon_j}\le1.
\]
On that simultaneous confidence event it implies every population error is at most its tolerance, by the triangle inequality. Vector/norm versions use corresponding valid simultaneous bounds. Training uncertainty must also be included when the claim averages over repeated training data. Evaluation on finitely many conditions does not establish a continuous-domain supremum without covering and regularity bounds.

\paragraph{Dynamics and well-posedness.}
For ODEs and SDEs give coefficient domains, initial laws and existence/uniqueness conditions in the intended solution sense, or keep the entire admitted solution set. For stochastic jump models specify nonnegative off-diagonal rates, diagonal normalization, rate orientation and non-explosion. For PDEs specify spatial domains, boundary/initial conditions, conservation or source terms, and a weak or classical solution class. Moving-domain concentration equations include dilution/transport terms. A stationary law is used only for a time-homogeneous augmented model under existence assumptions; a stationary approximation to a driven system must be justified. Absorbing extinction states are retained. Survival-conditioned, quasi-stationary and unconditional laws are distinct and must not be substituted for one another.

\paragraph{Identification and causal meaning.}
Two parameter values are observationally equivalent when they induce the same full distribution under all allowed observation/intervention designs, not merely the same likelihood value for one sample. A target is identified only if constant on the equivalence class. Gauge symmetries, label switching and reference conventions are quotiented or fixed. Finite-data compatible sets use a declared confidence, likelihood or discrepancy rule; an empty set signals model rejection. Bayesian credible sets are not automatically confidence sets. A prior or regularizer does not create data identification. Bounds intended to resolve a question must be informative at a predeclared width or decision threshold.

A causal contrast uses well-defined intervention versions and a stated causal model. Random assignment, or explicit consistency, overlap and exchangeability assumptions, is required for causal interpretation from observations. Interference and time-varying confounding are represented when present. Fixing a mediator by intervention differs from conditioning on its observed value. Marginal potential-outcome laws do not determine individual cross-world effects or an arbitrary joint coupling. Ethically admissible designs, rather than unconstrained interventions, define the experimental menu.

\paragraph{Metrics, units, optimization and limits.}
Specify the sampling or reference measure for densities and integrals, the log base, and dimensionless arguments of logarithms. Total variation is $\sup_A|P(A)-Q(A)|$; Wasserstein distances require the declared ground metric and appropriate moment conditions; KL divergence uses its stated absolute-continuity convention. All normed readouts are scaled consistently, and error probabilities use $0<\alpha<1$. A ratio has a positive denominator or an explicit undefined/zero convention. Thermodynamic entropy production is nonnegative in expectation under the relevant hypotheses; individual stochastic realizations need not be. An infimum is not an attained minimum without compactness/coercivity and continuity or another existence argument. Fix the order of large-population, long-time and vanishing-mutation limits. Stability claims declare perturbation space and symmetries.

\subsection*{How the provisional difficulty order was assigned}
Four editorial components are scored from 1 (lower burden) to 5 (higher burden): (a) formal or inferential difficulty, (b) access to evidence and identification burden, (c) mechanistic/computational integration, and (d) independent validation and transfer. Their unweighted total orders the entries. For selected entries in the highest bands, an explicit editorial priority breaks equal totals; all other ties are broken by original problem number. Equal totals are \emph{not} evidence of reliably different true difficulties. Scores are judgments, not measurements, completion-time estimates or consensus rankings.

The object ranked is an informative, nondegenerate resolution of the \emph{scoped} scientific target, with the stated budget and validation domain. A sharp useful impossibility result can qualify; an already known identity, vacuous bound, arbitrary favorable model or lookup table cannot. Rank 330 is merely the least difficult relative assessment in this catalogue, not a claim of an easy biological solution. Broad historical titles do not automatically outrank precisely formulated mechanistic tasks.


\section*{Selective source-check ledger}
\subsection*{SC09}\label{check:SC09}

Attwater et al., Trinucleotide substrates under pH-freeze-thaw cycles enable open-ended exponential RNA replication by a polymerase ribozyme (Nature Chemistry, 28 May 2025), DOI 10.1038/s41557-025-01830-y.

\textbf{Applied here to problems:} 284, 287.

\textbf{Access level:} Primary-paper abstract and selected HTML introduction/results, plus metadata.

\textbf{Support and limit:} Repeated RNA replication and partial ribozyme self-replication are important positive results. The stronger target additionally requires a connected chemically admissible pathway, hereditary renewal and the explicitly declared resource boundary.

\url{https://www.nature.com/articles/s41557-025-01830-y}

\url{https://pubmed.ncbi.nlm.nih.gov/40437192/}

\subsection*{SC10}\label{check:SC10}

Gianni et al., A small polymerase ribozyme that can synthesize itself and its complementary strand (Science 391, 1022-1028, 2026), DOI 10.1126/science.adt2760.

\textbf{Applied here to problems:} 284, 286, 287.

\textbf{Access level:} Indexed primary-paper abstract and publication metadata; direct full text not audited.

\textbf{Support and limit:} Copying the ribozyme and its complementary sequence is not identical to sustained autonomous population renewal with selectable variation under a complete prebiotic resource account.

\url{https://pubmed.ncbi.nlm.nih.gov/41678588/}

\url{https://www.science.org/doi/10.1126/science.adt2760}

\subsection*{SC11}\label{check:SC11}

Papastavrou, Horning and Joyce, RNA-catalyzed evolution of catalytic RNA (PNAS, 2024), DOI 10.1073/pnas.2321592121.

\textbf{Applied here to problems:} 286, 287.

\textbf{Access level:} Indexed primary-paper abstract and author-publication summary.

\textbf{Support and limit:} RNA-catalyzed replication and evolution provide partial positive results. They do not make the elementary single-peak threshold formula a universal theorem for arbitrary functional RNA networks.

\url{https://pubmed.ncbi.nlm.nih.gov/38437533/}


\end{document}
